	\section{Implementación y evaluación del prototipo}
	\label{cap:implementacion}
	Este capítulo documenta cómo se construyó, se desplegó y se verificó la plataforma diseñada en
	el Capítulo~\ref{sec:descripcion_solucion}. Se organiza siguiendo las capas del sistema: primero
	el entorno y las herramientas, después la implementación del lado del servidor, en seguida las
	aplicaciones web, luego el empaquetado en contenedores y la publicación en los servidores, y
	finalmente la evaluación del clasificador y las pruebas realizadas. Cierra con un balance
	explícito de lo que quedó implementado y lo que no, porque un documento técnico que solo enumera
	logros no permite a nadie retomar el trabajo.
	
	\subsection{Entorno de desarrollo y herramientas}
	\label{sec:entorno_herramientas}
	La plataforma se construyó con el conjunto de tecnologías cuyos fundamentos se expusieron en la
	Sección~\ref{sec:tecnologias}. La Tabla~\ref{tab:herramientas} resume la asignación concreta de
	cada herramienta a su función dentro del proyecto.
	
	\rowcolors{2}{rowblue}{white}
	
	\small
	\renewcommand{\arraystretch}{1.25}
	\setlength{\tabcolsep}{4pt}
	
	\begin{longtable}{|p{3.4cm}|p{3.2cm}|p{7.5cm}|}
		\caption{Herramientas empleadas y su función en el proyecto}
		\label{tab:herramientas} \\
		\hline
		\rowcolor{headerblue}
		\headercell{Capa} & \headercell{Herramienta} & \headercell{Función} \\
		\hline
		\endfirsthead
	
		\hline
		\rowcolor{headerblue}
		\headercell{Capa} & \headercell{Herramienta} & \headercell{Función} \\
		\hline
		\endhead
	
		Lenguaje del servidor & Java 21 & Lenguaje de los siete microservicios \\ \hline
		Marco del servidor & Spring Boot & Autoconfiguración, servidor embebido, seguridad y
		persistencia \\ \hline
		Construcción & Maven & Gestión de dependencias y empaquetado en archivo ejecutable \\ \hline
		Persistencia & Spring Data JPA & Mapeo objeto-relacional y derivación de consultas \\ \hline
		Base de datos & PostgreSQL 16 & Almacén relacional único compartido \\ \hline
		Autenticación & JSON Web Token & Tokens firmados con algoritmo simétrico \\ \hline
		Documentación de la API & OpenAPI 3 & Especificación e interfaz de exploración por servicio
		\\ \hline
		Hojas de cálculo & Apache POI & Generación del historial exportable y lectura de la
		importación masiva del catálogo \\ \hline
		Cliente web & Angular 21 & Las tres aplicaciones de página única \\ \hline
		Contenedores & Podman & Empaquetado y ejecución de los doce contenedores \\ \hline
		Servidor web & Nginx & Proxy inverso público y entrega de archivos estáticos \\ \hline
		Certificados & Certbot & Emisión y renovación automática del certificado de transporte
		\\ \hline
		Control de versiones & Git & Historial del proyecto \\ \hline
		Clasificador & Python, scikit-learn, spaCy, Sentence-BERT & Vectorización, entrenamiento y
		evaluación del componente de clasificación \\ \hline
	\end{longtable}
	
	\normalsize
	\setlength{\tabcolsep}{6pt}
	\renewcommand{\arraystretch}{1}
	
	Una característica del proyecto que conviene señalar desde el principio: cada microservicio es un
	proyecto de construcción independiente, sin un descriptor padre que los agrupe. La decisión
	simplifica la evolución de cada servicio por separado, pero tiene una consecuencia operativa
	concreta: no existe un comando único que compile todo el sistema, y cada servicio debe compilarse
	desde su propio directorio.
	
	\subsection{Implementación del lado del servidor}
	\label{sec:impl_backend}
	Los siete microservicios comparten deliberadamente la misma estructura interna. Esta uniformidad
	fue una decisión de implementación, no una consecuencia: cuando se construyó cada servicio nuevo
	se replicó la organización de carpetas, los filtros y el manejo de errores del anterior, en lugar
	de introducir un estilo distinto. El beneficio es que quien conoce un servicio conoce los siete, y
	un defecto corregido en uno se localiza sin dificultad en los demás.
	
	\subsubsection{Estructura común de un microservicio}
	Cada servicio organiza su código en cinco paquetes que corresponden a las capas descritas en la
	Sección~\ref{sec:arquitectura_software}:
	
	\begin{itemize}
		\item \texttt{controller}. Puntos de acceso HTTP. Reciben y devuelven objetos de
		transferencia, nunca entidades de persistencia, y declaran el rol requerido para invocarlos.
	
		\item \texttt{service}. Lógica de negocio y validación de las reglas. Es la única capa que
		conoce las reglas del dominio; los controladores no validan reglas y los repositorios no las
		conocen.
	
		\item \texttt{repository}. Acceso a datos mediante interfaces de persistencia que derivan las
		consultas a partir del nombre de sus métodos, sin escribir sentencias SQL.
	
		\item \texttt{entity}. Entidades de persistencia con su correspondencia a tablas y columnas.
	
		\item \texttt{model}. Objetos de transferencia de datos para las peticiones y las respuestas.
	\end{itemize}
	
	La separación entre \texttt{entity} y \texttt{model} merece una justificación explícita, porque
	implica escribir clases que a primera vista parecen redundantes. Exponer directamente la entidad
	de persistencia en una respuesta HTTP acopla el contrato público del servicio a la estructura de
	la base de datos: cualquier columna nueva aparece automáticamente en la respuesta, incluidas las
	que no deberían ser públicas ---la credencial cifrada de un usuario, por ejemplo---. Mantener
	objetos de transferencia separados obliga a decidir de forma consciente qué se expone.
	
	\subsubsection{Verificación de tokens y control de acceso}
	El mecanismo de autenticación se implementó como un filtro que se ejecuta antes de que la
	petición llegue al controlador. Su lógica es la misma en los siete servicios: extrae el token del
	encabezado de autorización, verifica su firma con la clave compartida, extrae de él la identidad
	y el rol, y establece el contexto de seguridad de la petición. Si el token falta, está vencido o
	su firma no coincide, el filtro rechaza la petición sin llegar a la lógica de negocio.
	
	Sobre ese contexto se aplica la autorización, declarada en cada controlador mediante anotaciones
	que exigen un rol determinado. Un detalle de configuración que resultó relevante durante la
	implementación: cuando un mismo prefijo de ruta tiene una parte pública y una parte
	administrativa, el orden de declaración de las reglas importa, porque el marco de seguridad
	aplica la primera coincidencia. Declarar la regla pública antes de la administrativa hace que la
	menos restrictiva prevalezca por ser más genérica, y el resultado es un punto de acceso
	administrativo efectivamente abierto. Esta situación se atendió declarando siempre la regla más
	específica primero, y quedó registrada como regla de negocio RN-A05 justamente para que no se
	reintroduzca al agregar servicios nuevos.
	
	\subsubsection{Manejo centralizado de errores}
	Cada servicio implementa un manejador global de excepciones que traduce cualquier error de la
	capa de negocio a una respuesta con estructura uniforme: mensaje, marca de tiempo y código. La
	alternativa ---capturar excepciones en cada controlador--- se descartó porque produce respuestas
	de error con formatos distintos según el punto de acceso, y obliga al cliente a implementar
	tantas variantes de manejo de error como puntos de acceso consuma. Con el manejador centralizado,
	las tres aplicaciones web leen el mensaje de error siempre del mismo lugar.
	
	\subsubsection{Documentación de la interfaz de programación}
	Los siete servicios exponen su especificación en formato OpenAPI y una interfaz de exploración
	interactiva. Esta decisión resolvió un problema concreto de coordinación: durante el desarrollo
	del cliente, la única forma de conocer la forma exacta de una petición era leer el código del
	controlador correspondiente. Con la especificación publicada, el contrato de cada servicio es
	consultable sin abrir el código.
	
	\subsubsection{Puntos de acceso implementados}
	La Tabla~\ref{tab:endpoints} presenta los puntos de acceso principales agrupados por servicio,
	con su condición de acceso. La especificación completa de cada uno, con sus parámetros y
	respuestas, se consulta en la interfaz de exploración de cada servicio y se reproduce en el anexo
	correspondiente.
	
	\rowcolors{2}{rowblue}{white}
	
	\footnotesize
	\renewcommand{\arraystretch}{1.2}
	\setlength{\tabcolsep}{4pt}
	
	\begin{longtable}{|p{2.6cm}|p{8.6cm}|p{3.0cm}|}
		\caption{Puntos de acceso principales por microservicio}
		\label{tab:endpoints} \\
		\hline
		\rowcolor{headerblue}
		\headercell{Servicio} & \headercell{Punto de acceso} & \headercell{Acceso} \\
		\hline
		\endfirsthead
	
		\hline
		\rowcolor{headerblue}
		\headercell{Servicio} & \headercell{Punto de acceso} & \headercell{Acceso} \\
		\hline
		\endhead
	
		\texttt{auth-service} & \texttt{POST /api/auth/login} & Público \\ \hline
		& \texttt{POST /api/auth/activar-cuenta} & Público \\ \hline
		& \texttt{POST /api/auth/solicitar-codigo} & Público \\ \hline
		& \texttt{POST /api/auth/reset-password} & Público \\ \hline
		& \texttt{GET} y \texttt{PUT /api/auth/me}, \texttt{/api/auth/perfil} & Sesión de quejoso
		\\ \hline
	
		\texttt{queja-service} & \texttt{POST /api/quejoso/\allowbreak quejas/registro-publico} & Público
		\\ \hline
		& \texttt{GET /api/quejoso/quejas/folio/\{folio\}} & Público \\ \hline
		& \texttt{POST /api/quejoso/quejas/registrar} & Sesión de quejoso \\ \hline
		& \texttt{GET /api/quejoso/quejas/mias} & Sesión de quejoso \\ \hline
		& \texttt{GET} y \texttt{PUT /api/quejoso/\allowbreak quejas/mias/\{folio\}} & Sesión de quejoso \\ \hline
		& \texttt{GET} y \texttt{PUT /api/quejoso/conciliaciones/**} & Sesión de quejoso \\ \hline
	
		\texttt{notificaciones-\allowbreak service} & \texttt{GET /api/notificaciones/mias} y
		\texttt{/mias/no-leidas} & Sesión de quejoso \\ \hline
		& \texttt{PUT /api/notificaciones/\{id\}/leida} & Sesión de quejoso \\ \hline
		& \texttt{POST /api/notificaciones/enviar} y \texttt{/registrar} & Uso interno entre
		servicios \\ \hline
	
		\texttt{catalogo-\allowbreak service} & \texttt{GET /api/catalogos/dependencias} y
		\texttt{/dependencias/\{clave\}} & Público \\ \hline
		& \texttt{GET}, \texttt{POST}, \texttt{PUT /api/catalogos/\allowbreak dependencias/admin/**} &
		Administración \\ \hline
		& \texttt{POST /api/catalogos/\allowbreak dependencias/admin/\allowbreak importar-excel} & Administración \\ \hline
	
		\texttt{admin-service} & \texttt{POST /api/admin/auth/login} & Público \\ \hline
		& \texttt{GET}, \texttt{POST}, \texttt{PUT}, \texttt{DELETE /api/admin/personal/**} &
		Administración \\ \hline
		& \texttt{GET} y \texttt{PUT /api/admin/plantillas/**} & Administración \\ \hline
		& \texttt{GET} y \texttt{POST /api/admin/\allowbreak seguridad/respaldos/**} & Administración \\ \hline
		& \texttt{POST /api/admin/seguridad/restaurar} & Administración \\ \hline
		& \texttt{GET /api/admin/seguridad/bitacora} & Administración \\ \hline
		& \texttt{GET /api/admin/dashboard/resumen} & Administración \\ \hline
		& \texttt{PUT /api/admin/perfil/password} & Cualquier cuenta de personal \\ \hline
	
		\texttt{revision-\allowbreak service} & \texttt{GET /api/revision/bandeja} & Recepción \\ \hline
		& \texttt{GET /api/revision/quejas/\{folio\}} y \texttt{/antecedentes} & Recepción \\ \hline
		& \texttt{GET /api/revision/\allowbreak quejas/evidencias/\{id\}} & Recepción \\ \hline
		& \texttt{POST /api/revision/\allowbreak quejas/\{folio\}/rechazar} y \texttt{/turnar} & Recepción
		\\ \hline
		& \texttt{POST /api/revision/registro-manual} & Recepción \\ \hline
		& \texttt{GET /api/revision/historial} y \texttt{/historial/exportar} & Recepción \\ \hline
		& \texttt{GET /api/revision/catalogos/areas} y \texttt{/defensores} & Recepción \\ \hline
		& \texttt{POST} y \texttt{GET /api/revision/conciliaciones} & Recepción \\ \hline
	
		\texttt{chatbot-\allowbreak service} & \texttt{GET /api/chatbot/menu} & Público \\ \hline
		& \texttt{GET}, \texttt{POST}, \texttt{PUT}, \texttt{DELETE /api/chatbot/\allowbreak admin/preguntas/**} &
		Administración \\ \hline
	\end{longtable}
	
	\normalsize
	\setlength{\tabcolsep}{6pt}
	\renewcommand{\arraystretch}{1}
	
	\subsection{Implementación de las aplicaciones web}
	\label{sec:impl_frontend}
	Las tres aplicaciones se desarrollaron como aplicaciones de página única con componentes
	autónomos, sin el sistema de módulos de versiones anteriores del marco de trabajo. Cada
	aplicación es un proyecto independiente con su propio árbol de rutas, su guarda de
	autenticación y sus interceptores.
	
	\subsubsection{Estructura común}
	Las tres comparten la misma organización: una capa \texttt{core} con los servicios de
	comunicación con el servidor, los modelos de datos y la configuración de seguridad del cliente;
	una capa \texttt{shared} con los componentes reutilizables ---encabezado, pie, avisos, ayuda
	flotante---; y una capa \texttt{pages} con un componente por pantalla, cargado de forma diferida
	para que el navegador no descargue el código de pantallas que el usuario no ha visitado.
	
	\subsubsection{Interceptores y guardas}
	Dos piezas transversales resuelven problemas que de otro modo habría que repetir en cada
	pantalla.
	
	El \textbf{interceptor de autorización} agrega el token de sesión al encabezado de cada petición
	saliente, de modo que ningún servicio de la aplicación tiene que recordar hacerlo. El mismo
	interceptor detecta las respuestas de sesión inválida o insuficiente y, en ese caso, cierra la
	sesión local y redirige al inicio de sesión. Esto implementa el requerimiento RNF-Q06 en un solo
	punto del código.
	
	La \textbf{guarda de rutas} impide el acceso a las secciones autenticadas cuando no existe una
	sesión válida en el cliente. Es importante entender su alcance: la guarda es una medida de
	usabilidad, no de seguridad, porque se ejecuta en el navegador y puede eludirse. La seguridad
	real la aplica el servidor, que rechaza cualquier petición sin token válido con independencia de
	lo que haga el cliente. Confundir ambas cosas es un error frecuente y por eso se hace explícito.
	
	\subsubsection{Detección de cambios sin zonas}
	La versión del marco de trabajo empleada permite ejecutar la aplicación sin la biblioteca que
	tradicionalmente intercepta las operaciones asíncronas para disparar la actualización de la
	vista. Ese modo reduce el trabajo del navegador, pero tiene una consecuencia que apareció durante
	la implementación: la vista no se actualizaba de forma automática al recibir la respuesta de una
	petición, de modo que los datos llegaban pero la pantalla seguía mostrando el estado anterior.
	
	La solución adoptada fue un interceptor adicional que solicita de forma explícita una
	actualización de la vista tras cada respuesta del servidor. Se documenta aquí porque es una
	particularidad no evidente del entorno y porque cualquier pantalla nueva que se agregue al
	proyecto depende de que ese interceptor siga registrado.
	
	\subsection{Empaquetado en contenedores y despliegue}
	\label{sec:impl_despliegue}
	El sistema se ejecuta en doce contenedores distribuidos en los dos servidores descritos en la
	Sección~\ref{sec:diseno_arquitectura}. La Tabla~\ref{tab:contenedores} presenta la matriz de
	despliegue.
	
	\rowcolors{2}{rowblue}{white}
	
	\small
	\renewcommand{\arraystretch}{1.2}
	\setlength{\tabcolsep}{4pt}
	
	\begin{longtable}{|p{3.4cm}|p{3.0cm}|p{7.6cm}|}
		\caption{Matriz de despliegue en contenedores}
		\label{tab:contenedores} \\
		\hline
		\rowcolor{headerblue}
		\headercell{Contenedor} & \headercell{Servidor} & \headercell{Contenido} \\
		\hline
		\endfirsthead
	
		\hline
		\rowcolor{headerblue}
		\headercell{Contenedor} & \headercell{Servidor} & \headercell{Contenido} \\
		\hline
		\endhead
	
		Base de datos & Aplicación & PostgreSQL 16 con almacenamiento persistente \\ \hline
		\texttt{auth-service} & Aplicación & Imagen dedicada del servicio de autenticación \\ \hline
		\texttt{quejas-service} & Aplicación & Imagen dedicada del servicio de quejas \\ \hline
		\texttt{notificaciones-\allowbreak service} & Aplicación & Imagen dedicada del servicio de
		notificaciones \\ \hline
		\texttt{catalogo-\allowbreak service} & Aplicación & Imagen dedicada del servicio de catálogo \\ \hline
		\texttt{admin-service} & Aplicación & Imagen propia que incluye el cliente de la base de
		datos, necesario para generar y restaurar respaldos \\ \hline
		\texttt{revision-\allowbreak service} & Aplicación & Imagen dedicada del servicio de revisión \\ \hline
		\texttt{chatbot-\allowbreak service} & Aplicación & Imagen dedicada del servicio de preguntas frecuentes
		\\ \hline
		Portal del quejoso & Acceso público & Servidor de estáticos con la aplicación compilada
		\\ \hline
		Consola de administración & Acceso público & Servidor de estáticos con la aplicación compilada
		\\ \hline
		Panel de revisión & Acceso público & Servidor de estáticos con la aplicación compilada
		\\ \hline
		Proxy inverso & Acceso público & Nginx con la configuración de enrutamiento y el certificado
		de transporte \\ \hline
	\end{longtable}
	
	\normalsize
	\setlength{\tabcolsep}{6pt}
	\renewcommand{\arraystretch}{1}
	
	\subsubsection{Configuración externalizada}
	Cada microservicio monta su configuración de producción como un volumen del contenedor, en lugar
	de incluirla en la imagen. La consecuencia práctica es una distinción importante en el
	procedimiento de operación: un cambio de configuración ---una dirección, un tiempo de
	expiración--- se aplica editando el archivo montado y reiniciando el contenedor, mientras que un
	cambio de lógica exige recompilar el servicio, generar un artefacto nuevo y reconstruir la
	imagen. Esta separación es la que permite que la misma imagen sirva en cualquier ambiente sin
	reconstruirse, tal como se expuso en la Sección~\ref{sec:tecnologias}.
	
	Durante la operación apareció una consecuencia no evidente de este esquema: editar el archivo de
	configuración en el servidor anfitrión no basta para que el contenedor lo tome, porque el punto de
	montaje conserva la referencia al archivo original. Es necesario reiniciar el contenedor, no solo
	recargar el servicio. Se documenta porque es una fuente de diagnósticos erróneos: el archivo se ve
	correcto en el servidor y el servicio sigue comportándose como antes.
	
	\subsubsection{Publicación segura y cifrado del canal}
	La publicación al exterior se resolvió en tres pasos. Primero se restringió el cortafuegos del
	servidor de aplicación para que los puertos de los microservicios y de la base de datos aceptaran
	conexiones únicamente desde la dirección del servidor de acceso público, dejando fuera cualquier
	otro origen de Internet. Después se configuró el proxy inverso del servidor de acceso público con
	las reglas de enrutamiento de la Tabla~\ref{tab:enrutamiento}. Por último se emitió el certificado
	de transporte mediante el método de validación por directorio web, se reconfiguró el proxy para
	atender el puerto cifrado y se dejó el puerto sin cifrar únicamente para redirigir el tráfico y
	atender la renovación automática del certificado.
	
	Durante este proceso se documentaron dos incidentes que conviene registrar por su valor
	diagnóstico. El primero fue una aparente imposibilidad de alcanzar los puertos web del servidor
	de acceso público desde el exterior, mientras que el puerto de administración remota respondía
	con normalidad. Tras descartar de forma sistemática la configuración del servicio, la publicación
	del contenedor y el cortafuegos local, se identificó que la causa era una regla de cortafuegos
	del proveedor que estaba guardada pero no sincronizada con la máquina virtual: existía en el panel
	de control pero no se había aplicado. El segundo fue un aviso recurrente del motor de contenedores
	sobre una red virtual mal configurada, que provenía de un intento previo de crear una red propia
	que finalmente no se utilizó; su eliminación resolvió el aviso. Ambos casos ilustran un principio
	de diagnóstico útil: cuando un componente responde de forma local y no desde el exterior, el
	problema no está en el componente sino en el camino.
	
	\subsection{Evaluación del clasificador de competencia}
	\label{sec:eval_clasificador}
	El clasificador se evaluó de forma independiente, conforme al protocolo definido en la
	Sección~\ref{sec:validacion_robusta}. Se construyeron ocho configuraciones que combinan las dos
	representaciones semánticas del texto con los cuatro clasificadores descritos en el marco
	teórico, y se evaluaron todas bajo el mismo procedimiento para que la comparación fuera válida.
	
	El protocolo aplicado tiene tres características que conviene explicitar, porque son las que
	impiden que las métricas resulten optimistas. La primera es la partición en tres subconjuntos:
	entrenamiento, validación cruzada dentro del entrenamiento, y un subconjunto de prueba aislado
	por completo durante el ajuste de hiperparámetros. La segunda es la estratificación de los
	pliegues, que preserva la proporción de clases en cada uno y evita que un pliegue quede sin
	ejemplos de la clase minoritaria. La tercera, y la más determinante, es que el criterio de
	selección del modelo final es su desempeño en modo autónomo ---sin los atributos jurídicos
	anotados--- porque ese es el escenario real de producción, en el que el quejoso redacta
	libremente sin conocer la normativa.
	
	% PENDIENTE: insertar aquí la tabla de resultados de las 8 configuraciones
	% (F1 en validación cruzada, desviación estándar entre pliegues, F1 en el subconjunto
	% de prueba con y sin atributos jurídicos) y las matrices de confusión del modelo
	% seleccionado. Falta que el autor proporcione las métricas.
	
	Los resultados numéricos de las ocho configuraciones, las matrices de confusión del modelo
	seleccionado y el análisis de los casos mal clasificados se presentan más adelante en esta misma
	sección.
	
	Una advertencia metodológica debe acompañar a cualquier lectura de estas métricas. Conforme a lo
	declarado en la Sección~\ref{sec:limitaciones}, el corpus con el que se entrenó y validó el
	clasificador es un conjunto sintético construido por el equipo de desarrollo, porque a la fecha
	de elaboración de este documento la Defensoría no ha entregado los expedientes históricos reales
	y anonimizados solicitados. En consecuencia, los valores reportados caracterizan el
	comportamiento del modelo sobre la distribución de ese corpus, no sobre la distribución real de
	las quejas que recibe la institución, y deben considerarse preliminares y sujetos a recalibración.
	
	\subsection{Pruebas y verificación}
	\label{sec:pruebas}
	La verificación del sistema se realizó en cuatro niveles, cada uno atendiendo un tipo distinto de
	defecto.
	
	\textbf{Verificación de compilación.} Cada microservicio se compila y empaqueta antes de
	desplegarse, y cada aplicación web se compila en modo de producción. Este nivel detecta errores
	de tipo y de plantilla, que en un lenguaje de tipado estático constituyen una fracción
	significativa de los defectos posibles.
	
	\textbf{Verificación de puntos de acceso.} Cada punto de acceso se probó de forma individual
	contra el servicio en ejecución, verificando el código de respuesta y la estructura del cuerpo.
	Este nivel detecta errores de configuración de seguridad ---un punto de acceso que debía ser
	público y exige token, o al contrario--- que la compilación no puede detectar.
	
	\textbf{Verificación de integración a través del proxy.} Cada ruta pública se probó desde el
	exterior, atravesando la cadena completa: proxy inverso, cortafuegos, microservicio y base de
	datos. Este nivel es el que detecta los problemas de enrutamiento y de cortafuegos, y fue el que
	permitió aislar los incidentes descritos en la Sección~\ref{sec:impl_despliegue}. El criterio de
	verificación aquí es específico: una respuesta de error generada por el servicio de destino
	constituye una prueba exitosa de conectividad, mientras que una expiración del tiempo de espera
	indica un problema en el camino.
	
	\textbf{Verificación funcional de los flujos.} Los flujos completos de cada actor se ejecutaron
	manualmente desde el navegador, siguiendo los casos de uso de la Sección~\ref{sec:casos_uso}:
	presentar una queja sin cuenta, activarla, iniciar sesión, editar, validar desde recepción,
	rechazar, turnar y consultar el historial exportado.
	
	Un límite de este esquema de verificación debe declararse con claridad: el proyecto no cuenta con
	pruebas automatizadas ---unitarias, de integración o de extremo a extremo--- ni con un flujo de
	integración continua que las ejecute. Toda la verificación descrita se realizó de forma manual.
	Esto significa que una modificación futura puede reintroducir un defecto ya corregido sin que
	nada lo detecte de forma automática, y constituye la deuda técnica más relevante del proyecto en
	materia de calidad.
	
	\subsection{Estado de avance por componente}
	\label{sec:estado_avance}
	La Tabla~\ref{tab:estado_avance} presenta el balance de lo implementado al cierre de esta etapa.
	Se incluye de forma deliberada la columna de lo no implementado, porque el valor de este balance
	para quien retome el proyecto está en lo que falta, no en lo que ya está.
	
	\rowcolors{2}{rowblue}{white}
	
	\small
	\renewcommand{\arraystretch}{1.25}
	\setlength{\tabcolsep}{4pt}
	
	\begin{longtable}{|p{3.7cm}|p{2.1cm}|p{8.2cm}|}
		\caption{Estado de avance por componente al cierre de la etapa}
		\label{tab:estado_avance} \\
		\hline
		\rowcolor{headerblue}
		\headercell{Componente} & \headercell{Estado} & \headercell{Observaciones} \\
		\hline
		\endfirsthead
	
		\hline
		\rowcolor{headerblue}
		\headercell{Componente} & \headercell{Estado} & \headercell{Observaciones} \\
		\hline
		\endhead
	
		Portal del quejoso & Operativo & Los diecisiete requerimientos funcionales del actor están
		implementados y desplegados \\ \hline
	
		Panel de revisión & Operativo & Los trece requerimientos del actor de recepción están
		implementados \\ \hline
	
		Consola de administración & Operativo & Los nueve requerimientos de administración están
		implementados \\ \hline
	
		Servicios de análisis, subdefensoría y resolución & No iniciado & Los tres roles existen como
		cuentas y como catálogo de turnado, pero no cuentan con interfaz ni permisos propios \\ \hline
	
		Integración del clasificador & Pendiente & El componente se desarrolló y evaluó de forma
		independiente; su exposición como servicio de la plataforma queda para la etapa siguiente
		\\ \hline
	
		Cierre de sesión por inactividad & Pendiente & La expiración del token limita la ventana de
		exposición, pero no equivale a un cierre por inactividad \\ \hline
	
		Pruebas automatizadas & No iniciado & Toda la verificación es manual; no existe integración
		continua \\ \hline
	
		Corrección de hallazgos de análisis estático & Pendiente & Las vulnerabilidades de
		neutralización de entradas y de generación pseudoaleatoria identificadas requieren remediación
		prioritaria \\ \hline
	
		Separación de los cortafuegos por servidor & Pendiente & Ambos servidores comparten el mismo
		grupo de reglas; no hay exposición real hoy, pero conviene separarlos \\ \hline
	\end{longtable}
	
	\normalsize
	\setlength{\tabcolsep}{6pt}
	\renewcommand{\arraystretch}{1}
